%
% Documento: Introdução
%

\chapter{Introdução}\label{chap:introducao}
\section{Objetivos}
\subsection{Objetivo geral}
A gestão de risco de derivativos constitui um dos pilares fundamentais das finanças quantitativas modernas, especialmente em ambientes de elevada incerteza e volatilidade (\textcite{cogley2008}). Os modelos clássicos, como o de Black-Scholes, pressupõem mercados completos, ausência de custos de transação e parâmetros fixos, como a volatilidade constante. Entretanto, tais condições raramente se verificam na prática, resultando em risco residual nas estratégias de hedge e comprometendo a eficiência da cobertura.

Nos chamados mercados incompletos, a impossibilidade de eliminar totalmente o risco decorre da presença de múltiplas fontes de incerteza, da variabilidade temporal dos parâmetros e de restrições operacionais que inviabilizam a replicação perfeita de carteiras (\textcite{fontanaStochasticControlPerspective2023}). Nesse contexto, o desafio central consiste em desenvolver estratégias de hedge que sejam adaptáveis às mudanças de regime e capazes de responder dinamicamente às oscilações de mercado.

A Engenharia de Controle e Automação oferece um conjunto consolidado de ferramentas matemáticas e computacionais para o tratamento de sistemas dinâmicos incertos, variantes no tempo e sujeitos a ruído — características análogas às que definem os mercados financeiros reais. Entre essas ferramentas, o controle adaptativo destaca-se pela capacidade de ajustar parâmetros do sistema em tempo real, de modo a minimizar o erro entre a resposta desejada e o comportamento observado.

Dessa forma, a aplicação de técnicas de controle adaptativo à gestão de risco de derivativos constitui uma abordagem inovadora e interdisciplinar, que permite a formulação de estratégias de hedge dinâmicas e robustas. Por meio da realimentação contínua do erro de hedge, torna-se possível ajustar o nível de exposição da carteira conforme as condições estocásticas do mercado evoluem, buscando a minimização do risco residual e o aumento da estabilidade financeira da posição.

O presente trabalho tem como objetivo desenvolver e avaliar um modelo de hedge adaptativo baseado em controle de realimentação, aplicável a mercados incompletos. Pretende-se comparar o desempenho dessa estratégia com abordagens tradicionais, utilizando métricas quantitativas de risco, como o erro quadrático médio e o Value at Risk (VaR), de modo a demonstrar a viabilidade técnica e a eficiência da aplicação de controle adaptativo em finanças quantitativas.


\section{Hipótese}\label{sec:hipotese}
A hipótese central deste trabalho é que a aplicação de técnicas de controle adaptativo em estratégias de hedge dinâmico de derivativos é capaz de reduzir o risco residual e melhorar a estabilidade da exposição da carteira frente à volatilidade variável e às imperfeições estruturais dos mercados incompletos.

Especificamente, considera-se que o uso de realimentação de erro e atualização adaptativa de parâmetros pode resultar em uma estratégia de hedge mais eficiente do que modelos estáticos ou baseados em suposições determinísticas, permitindo resposta mais rápida e robusta às flutuações de mercado.



\section{Justificativa}\label{sec:justificativa}
A crescente complexidade dos mercados financeiros, marcada pela volatilidade acentuada (\textcite{gaillardetzRiskControlStrategies2019}) e pela presença de incertezas estruturais, exige novas abordagens para a mitigação de risco. Os métodos convencionais de hedge, embora amplamente utilizados, partem de premissas idealizadas e, por isso, tendem a apresentar desempenho limitado em situações de instabilidade e não linearidade dos preços.

Nesse cenário, o controle adaptativo apresenta-se como uma alternativa tecnicamente consistente, capaz de integrar os princípios da Engenharia de Controle à realidade dinâmica dos mercados financeiros. A analogia entre um sistema físico sujeito a perturbações e um sistema financeiro exposto à volatilidade reforça a aplicabilidade dos conceitos de modelagem, realimentação e adaptação na gestão de risco de derivativos.

A justificativa para este estudo repousa, portanto, em dois eixos principais. O primeiro é o interesse científico, que reside na possibilidade de transferir metodologias consagradas da engenharia para o domínio das finanças quantitativas, contribuindo para a evolução de modelos híbridos entre as áreas. O segundo é o interesse prático, uma vez que a adoção de estratégias de hedge adaptativas pode proporcionar ganhos efetivos em robustez e eficiência para instituições financeiras, gestores de risco e plataformas de negociação automatizada.

Assim, o presente trabalho propõe-se a explorar e validar o potencial do controle adaptativo como ferramenta de apoio à tomada de decisão em contextos de incerteza, demonstrando a relevância e a aplicabilidade da Engenharia de Controle e Automação no enfrentamento de problemas complexos de risco financeiro.